\section{Additional Background}
\subsection{The creation of the ASB}
\label{appendix-a}
The creation of the AAFS Standards Board (ASB) was in large part motivated by the 2009 National Academies report, which recommended, amongst other things, the creation of standards for forensic science to ``ensure desirable characteristics of services and techniques such as quality, reliability, efficiency, and consistency among practitioners'' and ``maintain autonomy from vested interest groups''~\cite{national2009strengthening}. As opposed to other organizations that responded to this call, the ASB was specifically created in response to interest in the forensic science community, in developing an ANSI-accredited standards-developing organization~\cite{ambrosius2025celebrating}. As the organization describes on their website, ``In July 2015, AAFS received a \$1.5 million grant from the Laura and John Arnold Foundation over a four-year period to become an accredited [standards-developing organization] and provide the standards to the public free of charge. Since then, the [ASB] has formed 15 discipline-specific Consensus Bodies covering everything from DNA and toxicology to forensic nursing and wildlife forensics''~\cite{ambrosius2025celebrating}.

\section{Additional Methods}
\label{sec:appendix-b}
\subsection{Our interpretations of ASB 018's line-level requirements that relate to internal validation}
\label{sec:appendix-b-interpretations}
Below, for each ASB 018 line-level requirement that we deem relevant for internal validation studies (quoted in bold text), we provide our interpretation of the requirement. We understood two of these requirements to be tautological, or always be met, regardless of the internal validation study's specific practices (R4.1.3-d, R4.1.3-e). For several others, we interpreted the requirement to be met simply by the existence of the audit and audit report. For both of these groups of requirements, we explain our interpretation under the quoted requirement. For all other requirements included in our study, we phrased our interpretation as a question that we used to assess each audit report for compliance with the requirement, which we document below. In Appendix~\ref{appendix-b-compliance}, we document sections of each audit report that we take to fulfill each of our interpretations of this third group of line-level requirements.

In constructing our interpretations, we drew on ASB 018's own definitions of terms when available. Throughout, we aimed to reduce ambiguity while retaining as much of ASB 018's language as possible.


\textbf{Notation:} When a single numbered ASB 018 requirement contains multiple line-level requirements, we add letters after the ASB 018 requirement number to label each line-level requirement. For example, Requirement 4.1.3 contains five individual ``shall'' statements, which we distinguish using letters (e.g., R4.1.3a-e). Several of the line-level requirements we include in our analysis are pulled from Annex A of ASB 018, which we denote using ``annex'' (e.g., R4.1.3-annex).



\begin{itemize}
    \item \textbf{(R4.1) ``The laboratory shall validate a probabilistic genotyping system prior to its use for casework samples in the laboratory.''}~\cite[p.~3]{asb018} 
    \begin{itemize}
        \item Does the report describe the laboratory as having at least some part in conducting or reviewing the internal validation study?
    \end{itemize}
    
    \item \textbf{(R4.1.1-a) ``Validations shall include both developmental and internal studies.''}~\cite[p.~3]{asb018} 
    \begin{itemize}
        \item \textit{Line-level requirement simply requires the existence of an internal validation study.}
    \end{itemize}
    
    \item \textbf{(R4.1.1-b) ``Developmental validation shall not replace internal validation.''}~\cite[p.~3]{asb018}
    \begin{itemize}
        \item \textit{Line-level requirement simply requires the existence of an internal validation study.}
    \end{itemize}
    
    
    \item \textbf{(R4.1.3-a) ``Internal validation studies shall address the following: accuracy, sensitivity, specificity, and precision.''}~\cite[p.~3]{asb018}   
    \begin{itemize}
        \item (i) Does any part of the report address the software's accuracy (i.e., the software's ability to produce the expected likelihood ratio as calculated manually or with an alternate software program or application)?
        \item (ii) Does any part of the report address the software's sensitivity (i.e., the software's ability to produce inclusionary likelihood ratios for known contributors)?
        \item (iii) Does any part of the report address the software's specificity (i.e., the software's ability to produce exclusionary likelihood ratios for known non-contributors)?
        \item (iv) Does any part of the report address the software's precision (i.e., the variation in likelihood ratios calculated from repeated software analyses of the same input data using the same set of conditions/parameters)?
    \end{itemize}
    
    \item \textbf{(R4.1.3-b) ``These studies shall include internally generated case-type profiles of known composition that represent (in terms of number of contributors, mixture ratios, and total DNA template quantities) the range of actual casework samples intended for analysis with the system at the laboratory.''}~\cite[p.~3]{asb018}
    \begin{itemize}
        \item Does the report include internally generated samples of known composition that span a range of actual casework samples (in terms of the number of contributors, mixture ratios, and total DNA template quantities)?
    \end{itemize}

     \item \textbf{(R4.1.3-c) ``Studies shall not be limited to pristine DNA samples but shall also include compromised DNA samples (e.g., low template, degraded, and inhibited samples).''}~\cite[p.~3]{asb018}
    \begin{itemize}
        \item Does the report include any DNA samples that the lab either describes as ``compromised'' or ``challenging'', or is low template, degraded, or inhibited?
    \end{itemize}

    \item \textbf{(R4.1.3-d) ``The internal validation shall not exceed the scope of conditions tested in the developmental validation.''}~\cite[p.~3]{asb018}
    \begin{itemize}
        \item \textit{Because ASB 018's broad definition of developmental validation could also encompass any internal validation, including the internal validation at hand, this requirement is always met.}
    \end{itemize}

    \item \textbf{(R4.1.3-e) ``Case type profiles that fall outside the range of conditions explored in the developmental validation shall require additional developmental validation studies.''}~\cite[p.~3]{asb018}
    \begin{itemize}
        \item \textit{Condition triggering this requirement is never met given ASB 018's broad definition of developmental validation (see notes for R4.1.3-d).}
    \end{itemize}

    \item \textbf{(R4.1.3-annex) ``The laboratory shall perform sufficient studies to address the variability inherent to the various aspects of DNA testing, data generation, analysis and interpretation of data and user input parameters.''}~\cite[p.~5]{asb018}
    \begin{itemize}
        \item Does any part of the report test the software using a set of at least two samples that captures at least one source of variability in DNA profile generation, DNA profile interpretation by users, or analysis?
    \end{itemize}

    \item \textbf{(R4.1.4-a) ``Internal validation studies shall include evaluating user input parameters that vary run to run.''}~\cite[p.~3]{asb018}
    \begin{itemize}
        \item Does the report include any experiment that changes a user input?
    \end{itemize}

    \item \textbf{(R4.1.4-b) ``The effects of artifacts (e.g., stutter) and parameters that relate to the statistical algorithm (e.g., run time parameters for the software system that can vary from system to system) shall also be evaluated. [...] [T]he specific parameters to be tested shall be determined by the laboratory.''}~\cite[p.~3]{asb018}
    \begin{itemize}
        \item (i) Does any part of the report assess the effect of an artifact on some aspect of system behavior?
        \item (ii) Does any part of the report assess the effect of a parameter that relates to the statistical algorithm on some aspect of system behavior? Are these parameters determined by the laboratory?
    \end{itemize}


    \item \textbf{(R4.1.5-a) ``Internal validation studies shall also include the evaluation of multiple propositions for case type samples to aid in the development of propositions.''}~\cite[p.~3]{asb018}
    \begin{itemize}
        \item Does any part of the report use multiple propositions?
    \end{itemize}
    
    \item \textbf{(R4.1.5-b) ``Such studies shall also consider the effect of overestimating and underestimating the number of contributors.''}~\cite[p.~3]{asb018}
    \begin{itemize}
        \item Does any part of the report consider the effect of overestimating and underestimating the number of contributors on some aspect of system behavior?
    \end{itemize}

    \item \textbf{(R4.1.6) ``For internal validation, the laboratory shall evaluate both the appropriate sample types (i.e., number of contributors, mixture ratios, and template quantities) and the number of samples within each type to demonstrate the potential limitations and reliability of the software. The laboratory shall base this evaluation on the intended application of the software.''}~\cite[p.~3]{asb018}
    \begin{itemize}
        \item Does any part of the report claim that the study's evaluated samples are appropriate for the intended application of the software?
    \end{itemize}

    \item \textbf{(R4.4-annex) ``Additional validation or a performance check shall be based on the list of documented changes provided by the developer that accompany each updated version of the software installed in the laboratory.''}~\cite[p.~5]{asb018}
    \begin{itemize}
        \item If the report indicates that the lab has adopted a previous version of the software, does any part of the report suggest the study is based on information about documented changes provided by the developer?
    \end{itemize}

    \item \textbf{(R4.5) ``All validation and performance check studies conducted by the laboratory shall be documented and retained by the laboratory.''}~\cite[p.~4]{asb018}
    \begin{itemize}
        \item \textit{Because all of the audits we analyze in this study have been documented in the audit reports we are analyzing, this requirement is met for all audits in our study.}
    \end{itemize}
\end{itemize}



\subsection{Compliance assessment}
In Table~\ref{appendix-table-compliance}, we document sections of each audit report that we take to fulfill each of our interpretations of ASB 018's line-level requirements (see Appendix~\ref{sec:appendix-b-interpretations} for our interpretations). 

\textbf{Additional note on our assessment of PBSO internal validation study of STRmix v2.6.2 against R4.4-annex:} While not explicitly describing differences between v2.6.2 and the previously validated v2.4, the PBSO v2.6.2 report does describe developer-provided documentation accompanying v2.6.2 as a guide for designing the PBSO validation of v2.6.2 and deliberately avoids re-validating certain software behaviors it deemed irrelevant to the study.

\label{appendix-b-compliance}
\begin{table*}[h]
\begin{tabular}{|  l  |l  |l|l  |l|l  |}
\hline
 \textbf{Line-level Req.} & \textbf{OCME v2.4 (2016)} &  \textbf{CBI v2.5 (2018)} &\textbf{PBSO v2.6.2 (2019)} &  \textbf{OCME v2.7 (2021)} &\textbf{MSP v2.9.1 (2024)}\\ 
 \hline\hline
 4.1 & Introduction (p. 2) &  Introduction (p. 2) &Introduction (p. 2) &  Introduction (p. 1) &Introduction (p. 3) \\  
 \hline
 4.1.3-a (i) & Exp. 2 (p. 3-5) &  Sec. A (p. 4-5) &Sec. A (p. 5-7) &  Exp. 2 (p. 4-6) &Sec. A (p. 5-7) \\
 \hline
 4.1.3-a (ii) & Exp. 4 (p. 7-17) &  Sec. D (p. 7-32) &Sec. D (p. 8-15) &  Exp. 4 (p. 14-25) &Sec. D (p. 13-29)  \\
 \hline
 4.1.3-a (iii) & Exp. 4 (p. 7-17) &  Sec. D (p. 7-32) &Sec. D (p. 8-15) &  Exp. 4 (p. 14-25) &Sec. D (p. 13-29)  \\
 \hline
 4.1.3-a (iv) & Exp. 6 (p. 19-22) &  Sec. M (p. 52-55) &Sec. M (p. 25-31) &  Exp. 6 (p. 28-30) &Sec. M (p. 64-66)\\
 \hline
 4.1.3-b & (p. 2, 40) &  (p. 8, 56)  &(p. 2, 31) &  (p. 2, 64) &(p. 3, 14)  \\
 \hline
 4.1.3-c & Exp. 14 (p. 38-40) &  Sec. L (p. 47-52) &Sec. L (p. 25) &  Exp. 13 (p. 54-63) &Sec. J (p. 48-58)  \\
 \hline
 4.1.3-annex & Exp. 9 (p. 24-27) &  Sec. M (p. 52-53) &Sec. F (p. 18-22) &  Exp. 6 (p. 28-30) &Sec. M (p. 64-66)  \\
 \hline
 4.1.4-a & Exp. 10 (p. 27-33) &  Sec. F (p. 35-38) &Sec. F (p. 18-22) &  Exp. 9 (p. 36-44) &Sec. F (p. 30-40) \\
 \hline
 4.1.4-b (i) & Exp. 11 (p. 33-35) &  Sec. G (p. 40) &Sec. G (p. 22) &  Exp. 10 (p. 44-48) &Sec. G (p. 41-44)  \\
 \hline
 4.1.4-b (ii) & Exp. 6 (p. 19-22) &  Sec. M (p. 53-55) &Appendix 3 (p. 36-45) &  Exp. 6 (p. 28-30) &Sec. M (p. 64-66) \\
 \hline
 4.1.5-a & Exp. 5 (p. 18-19) &  Sec. E (p. 32-34) &Sec. E (p. 16-17) &  Exp. 5 (p. 25-27) &Sec. E (p. 29-30)  \\
 \hline
 4.1.5-b & Exp. 10 (p. 27-33) &  Sec. F (p. 35-38) &Sec. F (p. 18-22) &  Exp. 9 (p. 36-44) &Sec. F (p. 30-40)  \\
 \hline
 4.1.6 & Conclusion (p. 40) &  Conclusion (p. 56) &Conclusion (p. 31) &  Conclusion (p. 64) &Conclusion (p. 96)  \\
 \hline
  4.4-annex & N/A &  N/A &Introduction (p. 2) &  Introduction (p. 1-2) &N/A\\
 \hline
 
 
\end{tabular}
\caption{For each internal validation report, we document the page numbers of the internal validation report that contain text that satisfies our interpretation of each line-level requirement. When multiple sections of a report could satisfy our interpretation of a given line-level requirement, we document the page numbers for one of the sections. All page numbers refer to the numbered pages of the internal validation report. The OCME v2.4 audit report was updated in 2019 to correct transcriptional errors and improve the clarity of tables, figures and end notes---the original audit was completed in 2016.}
\label{appendix-table-compliance}
\end{table*}